\chapter{Thành và Bại}
\markboth{Thành và Bại}{Success and Failure}

\section{Một người thành công sớm nên tưởng mình luôn đúng.}
\begin{truyen}{Một người thành công sớm nên tưởng mình luôn đúng.}{Early success can create late blindness.}
\chuhoa{V}{ài cú thắng sớm làm anh tưởng mình nhìn đâu cũng chuẩn. Thế là bắt đầu bỏ ngoài tai lời can, xem nhẹ rủi ro, nhìn người khác bằng nửa con mắt. Thành công mà không kéo theo tỉnh táo thì khác gì tự bơm hơi cho cú ngã sau này.}
\textit{(A few early victories made him believe that his instincts were always right. So he stopped listening to warning, underestimated risk, and looked down on others. Success without awareness can become a slow kind of failure.)}
\end{truyen}
\begin{baihoc}
Thắng vài trận chưa nói lên gì. Nhiều người hỏng cũng từ lúc tưởng mình không thể hỏng.
\textit{(Success does not always make people wiser; sometimes it makes them more careless.)}
\end{baihoc}
\begin{ghinhoanh}
\textbf{Short English to remember:}

Success can blur judgment.
\end{ghinhoanh}
\begin{tuhocnhanh}
1. Khi làm tốt một việc, em có dễ chủ quan không? \textit{(When you do well, do you become overconfident?)}

2. Thành công nào của em cần đi kèm nhiều tỉnh táo hơn? \textit{(Which success of yours needs more awareness around it?)}
\end{tuhocnhanh}
\ngancach

\section{Một lần bại làm lộ ra người thật bên cạnh mình.}
\begin{truyen}{Một lần bại làm lộ ra người thật bên cạnh mình.}{Failure reveals the map of loyalty.}
\chuhoa{K}{hi đang lên, quanh anh đông như chợ. Đến lúc công việc gãy, người biến mất nhanh hơn tin nhắn seen. Chỉ còn vài người ở lại hỏi một câu tử tế. Thất bại đau thì đau, nhưng ít ra nó dọn bớt đám người chỉ đứng cạnh mình lúc có lợi.}
\textit{(When life was going well, he was always surrounded by noise and company. When work collapsed, only a few remained to ask, care, and support. Failure hurts, but it also wipes off the dust that makes us misread human loyalty.)}
\end{truyen}
\begin{baihoc}
Thua một cú mà nhìn ra ai thật ai giả thì cũng không hẳn lỗ.
\textit{(Failure does not only take away; it also gives a clearer map of who and what is real.)}
\end{baihoc}
\begin{ghinhoanh}
\textbf{Short English to remember:}

Failure reveals true company.
\end{ghinhoanh}
\begin{tuhocnhanh}
1. Khi em khó khăn, ai là người còn ở lại? \textit{(When you struggle, who still stays?)}

2. Em đã học được gì về các mối quan hệ từ một lần thất bại? \textit{(What have you learned about relationships from a failure?)}
\end{tuhocnhanh}
\ngancach

\section{Một học sinh bại ở lần đầu nhưng thành ở cách đứng dậy.}
\begin{truyen}{Một học sinh bại ở lần đầu nhưng thành ở cách đứng dậy.}{Rising well is also a success.}
\chuhoa{E}{m thi hỏng, chán tới mức muốn mặc kệ luôn. Nhưng chán xong vẫn quay lại hỏi thầy chỗ yếu, sửa cách học, cắm cúi làm lại từ đầu. Không phải ai ngã cũng thua; nhiều người chỉ thua thật khi nằm lì ở đó và tự thương thân mãi.}
\textit{(The student failed an exam and for a while wanted to abandon everything. But then came a return, a request for guidance, a change of method, and steady effort day after day. Sometimes what deserves to be called success is not the absence of falling, but the refusal to remain there.)}
\end{truyen}
\begin{baihoc}
Ngã chưa đáng sợ. Nằm luôn mới đáng sợ.
\textit{(Falling does not make us truly defeated; staying down too long does.)}
\end{baihoc}
\begin{ghinhoanh}
\textbf{Short English to remember:}

Getting up well is success.
\end{ghinhoanh}
\begin{tuhocnhanh}
1. Sau thất bại gần nhất, em đứng dậy theo cách nào? \textit{(After your most recent failure, how did you rise?)}

2. Em còn bước nào chưa làm để quay lại tốt hơn? \textit{(What step have you not yet taken to return better?)}
\end{tuhocnhanh}
\ngancach

\section{Thành công ngoài đời nhưng thất bại trong gia đình.}
\begin{truyen}{Thành công ngoài đời nhưng thất bại trong gia đình.}{Outer success can hide inner ruin.}
\chuhoa{A}{nh ngoài xã hội thì oai, hợp đồng ký đều, đi đâu cũng được nể. Nhưng về nhà, con ngại nói chuyện, vợ quen cảnh tự xoay, cha mẹ già muốn gặp cũng phải chờ. Có những người thắng ngoài đường mà thua sạch trong nhà.}
\textit{(He was respected at work, signed major deals, and traveled widely. Yet at home his child no longer wanted to talk, his wife had grown used to being alone, and his aging parents were rarely seen. Some successes shine brightly in the world while leaving home in darkness.)}
\end{truyen}
\begin{baihoc}
Nếu càng thành đạt mà nhà càng lạnh thì nên coi lại mình đang thắng cái gì.
\textit{(If we succeed outside while failing where love is closest, we need to revisit our definition of success.)}
\end{baihoc}
\begin{ghinhoanh}
\textbf{Short English to remember:}

Success outside can fail at home.
\end{ghinhoanh}
\begin{tuhocnhanh}
1. Điều em gọi là thành công có đang bỏ quên ai không? \textit{(Is your idea of success neglecting someone important?)}

2. Thành công nào đáng giá nếu phải đổi bằng những mối quan hệ cốt lõi? \textit{(What kind of success is worth the price of core relationships?)}
\end{tuhocnhanh}
\ngancach

\section{Bại một kế hoạch để thành một con người chín hơn.}
\begin{truyen}{Bại một kế hoạch để thành một con người chín hơn.}{A broken plan can mature a person.}
\chuhoa{C}{ô từng tưởng dự án này gãy là mình tiêu luôn. Đến lúc nó gãy thật, cô mới phải học xin lỗi, học nghe góp ý, học làm lại mà bớt đổ thừa. Kế hoạch hỏng, nhưng cái đầu bớt non đi một khúc.}
\textit{(She once thought that if this project failed, she would be finished. Yet when it failed, she was forced to learn how to apologize, receive feedback, and begin again without resentment. The plan failed, but the person inside her became more mature.)}
\end{truyen}
\begin{baihoc}
Có cái bại đập nát kế hoạch, nhưng nhờ vậy mới đập bớt ảo tưởng trong đầu.
\textit{(Some failures kill a plan but save a person from shallowness.)}
\end{baihoc}
\begin{ghinhoanh}
\textbf{Short English to remember:}

A broken plan can build a person.
\end{ghinhoanh}
\begin{tuhocnhanh}
1. Thất bại nào đã làm em chín lên nhiều nhất? \textit{(Which failure has matured you the most?)}

2. Em đang tiếc một kế hoạch hay đang lớn lên từ nó? \textit{(Are you only grieving a broken plan, or are you growing through it?)}
\end{tuhocnhanh}
\ngancach
